\documentclass[../structural_economics_book.tex]{subfiles}
\begin{document}
\ifSubfilesClassLoaded{\setcounter{chapter}{27}}{}
\chapter{The General Case}\label{ch:28}
\chaptersubtitle{Where the Standard Model Is the Special Case, and Structure Is Everything Else}
\begin{quote}
\textit{``O subdesenvolvimento n\~{a}o constitui uma etapa necess\'{a}ria do processo de forma\c{c}\~{a}o das economias capitalistas modernas. \'{E}, em si, um processo particular, resultante da penetra\c{c}\~{a}o de empresas capitalistas modernas em estruturas arcaicas.''}\\[0.3em]
\textnormal{[``Underdevelopment does not constitute a necessary stage in the formation of modern capitalist economies. It is, in itself, a particular process, resulting from the penetration of modern capitalist enterprises into archaic structures.'']}\\
\hfill --- Celso Furtado, \textit{Desenvolvimento e Subdesenvolvimento} (1961)
% TODO-VERIFY: this wording may come from Teoria e Politica do Desenvolvimento Economico (1967) rather than the 1961 book; the canonical 1961 line is "um processo historico autonomo, e nao uma etapa ...". Check both editions and give a page.
\end{quote}
\nocite{furtado1961desenvolvimento}

\vspace{1em}

This is the last chapter. It has three purposes: to state the book's thesis in its final, corrected form; to restate each chapter's result in terms of $W$ and the primitives it adds; and to turn what the book has not done into a research programme.


%% =====================================================================
\section{The Thesis}
\label{ch28:sec:thesis}
%% =====================================================================

\begin{remark}[Framing Principle: The General Case]\label{ch28:thm:general}
The benchmark competitive model abstracts from social structure. Agents interact anonymously, through prices: aggregate demand weights each agent by her expenditure share $s_i$, so the implied interaction matrix is rank one, $W^{\text{mkt}} = \ones\mathbf{s}'$, with $\Gamma = 1$ and $\boldsymbol{\pi} = \mathbf{s}$. With equal endowments it is $\Wce = \frac{1}{N}\ones\ones'$. Utility interdependence is absent ($W^U = 0$, no externalities), which is what the First Welfare Theorem needs; it holds for \emph{any} wealth distribution $\mathbf{s}$ (Chapter~1, Section~\ref{ch01:sec:conventions}).

Structural economics studies the general case: non-anonymous, row-stochastic interaction matrices $W$ --- production, information and utility networks --- and how they are produced and changed. Every $W$ can be written as
\begin{equation}\label{ch28:eq:departure}
W = \Wce + \Delta(W), \qquad \Delta(W) := W - \tfrac{1}{N}\ones\ones',
\end{equation}
and the departure is summarized by two coordinates that are independent of each other (Proposition~\ref{ch06:prop:gap-power-independent}):
\begin{itemize}
\item \textbf{power distance} $d_\pi = \tfrac12\|\boldsymbol{\pi} - \tfrac{1}{N}\ones\|_1$, or equivalently $\Neff/N$, which measures how unequally the structure distributes power;
\item \textbf{integration deficit} $1 - \Gamma$, which measures how slowly shocks, beliefs and demand propagate across the network.
\end{itemize}
Neither determines outcomes alone. Allocations, beliefs, preferences and welfare depend on $W$ \emph{and} on primitives that are not functions of $\Delta(W)$: the bargaining surplus $\Phi$ and curvature $\kappa$, the formation parameters $\rho$ and $\mu$, endowments, the technical coefficients $A$, the coalition game, the product space.
\end{remark}

This is a framing principle, not a theorem: the decomposition~\eqref{ch28:eq:departure} holds for every matrix, so its content lies entirely in the maps from $W$ and the primitives to outcomes. Those maps are what the previous 27 chapters study --- some with theorems, more with conjectures and illustrative calculations.


%% =====================================================================
\section{The Book in One Page}
\label{ch28:sec:corollaries}
%% =====================================================================

\label{ch28:prop:book}
The table restates each chapter's main result in terms of $W$ and lists the primitives, beyond $W$, on which it depends. Integration is $\Gamma$; power is $\boldsymbol{\pi}$, with concentration measured by $\Neff$, $\Var(\boldsymbol{\pi})$ or $G(\boldsymbol{\pi})$. Bolivian production-network, complexity and distribution values are sourced (Section~\ref{ch28:sec:bolivia}); $\sigma_B$ and the $W^I$ and $W^U$ values remain illustrative.

{\small
\begin{longtable}{@{}r>{\raggedright\arraybackslash}p{2.5cm}>{\raggedright\arraybackslash}p{7.6cm}>{\raggedright\arraybackslash}p{3.1cm}@{}}
\toprule
\textbf{Ch.} & \textbf{Topic} & \textbf{Result in terms of $W$} & \textbf{Primitives $X$} \\
\midrule
\endfirsthead
\toprule
\textbf{Ch.} & \textbf{Topic} & \textbf{Result in terms of $W$} & \textbf{Primitives $X$} \\
\midrule
\endhead
1 & Matrix notation & Allocation, production and interaction as matrices. The benchmark is anonymous market interaction $W^{\text{mkt}} = \ones\mathbf{s}'$; $\Wce$ is its equal-endowment case. & Endowments, preferences \\
2 & Welfare theorems & FWT holds for any $\mathbf{s}$ given complete markets, price-taking and no externalities ($W^U = 0$). Unequal $\boldsymbol{\pi}$ is distributional: $L_{\text{power}} = (\kappa N/2)\Phi^2\Var(\boldsymbol{\pi})$. & $\kappa$, $\Phi$, $U$ \\
3 & Growth \& distribution & Growth spreads across blocks at a speed set by the integration of $W^P$; multipliers come from $A$, not from $W$. & $A$, saving, technology \\
4 & Benchmark assumptions & Six assumptions restated as restrictions on $W$: $W^U = 0$, $W^P$ exogenous, no $W^I$, anonymity, $\boldsymbol{\pi} = \frac{1}{N}\ones$ (equal endowments), fixed $\theta$. & --- \\
5 & Three networks & $W^P$, $W^I$, $W^U$; a composite $\sum_k \alpha_k W^k$ exists only on a common node set. & $\alpha_k$ \\
6 & Power & $\boldsymbol{\pi}' W = \boldsymbol{\pi}'$; $\boldsymbol{\pi}$ uniform iff $W$ doubly stochastic; $\Gamma$ and $\boldsymbol{\pi}$ independent. & Reading of $\boldsymbol{\pi}$ \\
7 & Bargaining & With $\boldsymbol{\pi}$ as Nash weights (an assumption): $x_i = \Omega/N + (\pi_i - 1/N)\Phi$. & $d_i$, $S$, $\Phi$ \\
8 & Beliefs & DeGroot consensus $b^* = \boldsymbol{\pi}'\bb^{(0)}$, reached at a speed set by $\Gamma(W^I)$; $\Neff = 1/\sum(\pi_j^I)^2$. & $\bb^{(0)}$ \\
9 & Preferences & $\theta_{i,t+1} = \rho\theta_{i,t} + (1-\rho)\theta^{\circ} + \mu(\sum_j W^I_{ij}r_j - r^*)$: preferences formed through $W^I$. Definition: $\Psi_i(\Delta W, T) = F_i/E_i$, for a stated perturbation and horizon. & $\rho$, $\mu$, $\theta^{\circ}$, $r^*$ \\
10 & Formation model & The loop $\boldsymbol{\theta} \to \x^* \to \mathbf{r} \to \boldsymbol{\theta}$ is a map $T$ with Jacobian $J = \rho I + \mu\phi W^I C P$; unique stable steady state if $\rho(J) < 1$ (affine case); traps need nonlinear returns. & $\rho$, $\mu$, returns \\
11 & Measuring $\Psi$ & The cohort--year ratio $\hat{\Psi}_{\text{var}}$ is descriptive, not $\Psi$; identifying $\Psi$ needs strong assumptions or randomized $W^I$. & Identification design \\
12 & Welfare & $W^I$ shapes the preferences that evaluate $W^P$ (``the ruler is rubber''). Under fixed aspirations, structures rank by $\Var(\boldsymbol{\pi})$, not by $\Gamma$. & Welfare criterion \\
13 & Network evolution & $W_{t+1} = \mathcal{E}(W_t; \x^*_t)$: proximity-driven edge creation and destruction; structural equilibria exist and can be multiple (illustratively $\Gamma = 0.072$ and $0.643$, same $\boldsymbol{\pi}$). & Creation, destruction rates \\
14 & Coalitions & Payoff $v_i \approx (1+\Psi)(\Delta\Omega/N + \Phi\,\Delta\pi_i)$: the power part is constant-sum; alliance cycling is a conjecture. & Game form \\
15 & Operator design & Choose $\Delta W$ to raise integration and to lower concentration; because $\Gamma \perp \boldsymbol{\pi}$, raising $\Gamma$ need cost no one power. & Budget, costs \\
16 & Global $W^P$ & Global input and trade shares; ECI comes from the country $\times$ product matrix $M$, not from $W$. & $M$, RCA rule \\
17 & Domestic diagnostics & $\Gamma$, $\sigma_B$, $G(\boldsymbol{\pi})$, $\Neff$, $\Psi$, read separately. Bolivia 2019: $\Gamma = 0.467$ (0.475 like-for-like, inside the comparator range 0.463--0.649; 0.33--0.69 across normalizations), $G(\boldsymbol{\pi}^P) = 0.445$ (Korea 0.658). The enclave shows in $\pi_{\text{hydrocarbons}} = 0.039$ with 70\% of output exported, not in $\Gamma$. Illustrative: $\sigma_B = 0.62$, $\Neff^I = 3.4$. & Data, aggregation, normalization \\
18 & Gender \& race & Node attributes constrain $W$ through assortative mixing $\lambda_k$, which lowers $\Gamma$ when within-group links dominate; power gaps need unequal group size or opportunity structure. Intersectionality is the interaction term. & $\lambda_k$ \\
19 & Scale invariance & $\pi_I = \pi_G \cdot \pi_{N|c} \cdot \pi_{\text{within}}$ exactly (stochastic complementation); under exact lumpability aggregation raises $\Gamma$; isomorphism is a conjecture. & Nesting of scales \\
20 & SAMs & A row-normalized SAM is a $W$ whose Perron vector is account size; multipliers come from the SAM coefficient matrix. & SAM accounts \\
21 & $W^I$ measurement & Media overlap, reference groups, social networks, digital traces; audience data measure integration and reach, not who is believed. Bolivia: $\Gamma^I = 0.74$, $\Neff = 5.59$ of 6 groups; no information bubbles in audience shares. & Survey design \\
22 & $W^U$ measurement & Transfers, experiments, voting, surveys. Bolivia: assortative ($\Gamma = 0.30$) but not concentrated ($\Neff^U = 3.85$ of 4 groups). & Survey design \\
23 & Time series & Per-period $\Gamma_t$ does not control $W_t\cdots W_s$: use $\tau_1$, the realized-product gap and the absolute probability sequence. Bolivia's $\Gamma$ fell from 0.567 (1990) to 0.484 (2013), $v_\Gamma = -0.036$ per decade; Korea's fell from 0.626 (1995) to 0.511 (2019). Aggregate $\Gamma$ does not separate development paths, and the SCI is undefined when $\Gamma$ starts above $\Gamma^{\text{ref}}$, as both do. & Filter, $\Gamma^{\text{ref}}$ \\
24 & Hegemony & Hegemonic equilibrium: $W^I$-formed consent for $W^P$. Under conformity, counter-hegemony succeeds past a saddle-node threshold; the cycle is a conjecture. & $\rho$, $\beta$, $s$ \\
25 & Truncation & Removing a bridge set bounds $\Gamma$ by twice the remaining conductance; that the capital-goods core is a bridge is a conjecture. It gets no support at 12 sectors: Bolivia's $\Gamma$ (0.475) is close to Korea's (0.511) while its ECI ($-1.30$) is far below ($+1.81$). Testing it needs 39-industry or product-level data. & $\phi$, resolution \\
26 & Synthesis & Conjecture: sustained growth with equity needs $\Gamma > \Gamma^*$, ECI $> 0$ and $\Neff > N^*_{\text{eff}}$. Bolivia's 2005--2019 fall in inequality (Gini 0.600 to 0.426) with a negative, falling ECI contradicts it as a claim about equity gains. & Thresholds \\
27 & Algorithms & Engagement ranking lowers $\lambda_k$ and $\Gamma$ on the sorted attribute; viral hubs lower $\Neff$ but raise $\Gamma$. & Engagement $e_{ij}$ \\
\bottomrule
\end{longtable}
}


%% =====================================================================
\section{What Furtado Knew}
\label{ch28:sec:furtado}
%% =====================================================================

Celso Furtado's \citeyear{furtado1961desenvolvimento} thesis --- that underdevelopment is not a stage but a structure --- is the sentence that generated this book. Every chapter is an elaboration.

\begin{remark}[Furtado's Thesis in the Language of $W$]\label{ch28:prop:furtado}
\citet{furtado1961desenvolvimento} made three claims. Their translation is ours, and it is a reinterpretation rather than a proof:
\begin{enumerate}
\item \textbf{Underdevelopment is not the absence of development.} It is a specific structural configuration, the product of the \emph{interaction} of a modern sector, tied to the world economy, with an archaic one. In $W$ terms it has \emph{two} features that the book's two coordinates keep apart. The first is an \emph{enclave}: a modern export sector weakly tied to the domestic production network. It shows in the sector's small Perron weight in $W^P$ and its high export share, not necessarily in a small aggregate $\Gamma(W^P)$. The second is an extraction channel through which surplus flows from the archaic to the modern sector, which concentrates bargaining power ($\boldsymbol{\pi}$ far from uniform under Assumption~7.A, Chapter~7). Neither implies the other. In Bolivia's measured $W^P$ the enclave is there (hydrocarbons: $\pi = 0.039$, 70\% of output exported), while aggregate $\Gamma$ is in the comparator range and production power is less concentrated than Korea's.

\item \textbf{The structure reproduces itself.} Underdevelopment is a stable fixed point, not a transient state. The mechanisms of reproduction are: proximity-driven edge formation (Chapter~13), self-confirming preferences (Chapter~10), coalition dynamics that protect the incumbent $W$ (Chapter~14), and hegemonic formation that produces consent (Chapter~24).

\item \textbf{The structure can be changed, but not by the standard tools.} Price adjustments, trade liberalization, and transfers do not change $W$. They change $\boldsymbol{\omega}$ (endowments) and relative prices within the existing structure. Changing $W$ requires the operators of Chapter~15: information policy, institutional reform, education, infrastructure, and industrial policy, sequenced by $\Psi$ and constrained by the political economy of Chapter~14.
\end{enumerate}

Every chapter in this book develops one of these three claims:
\begin{center}
\begin{tabularx}{\textwidth}{LLL}
\toprule
\textbf{Furtado's claim} & \textbf{Formal content} & \textbf{Chapters} \\
\midrule
Not a stage but a structure & Non-anonymous $W$: an export enclave (low $\pi$, high export share) and concentrated bargaining power & 1--4 \\
Self-reproducing & Stable fixed point of the co-evolutionary system & 5--12, 24 \\
Changeable but not by standard tools & Operators on $W$, not on $\boldsymbol{\omega}$ & 13--15, 25--26 \\
\midrule
Measurement & SAMs, $W^I$, $W^U$, trajectories & 16--23 \\
Scale invariance & Same at all scales & 19 \\
Contemporary & Algorithmic $W^I$ & 27 \\
\bottomrule
\end{tabularx}
\end{center}
\end{remark}

\begin{intuition}[Why ``the general case'']
The benchmark Arrow--Debreu model abstracts from social structure: agents interact anonymously through prices ($W^{\text{mkt}} = \ones\mathbf{s}'$), there are no externalities ($W^U = 0$), and preferences are taken as given rather than formed through $W^I$. This is a deliberate abstraction, not a mistake, and modern economics has relaxed each part of it --- production-network macroeconomics \citep{long1983real,acemoglu2012network,baqaee2019macroeconomic}, network learning \citep{degroot1974reaching,golub2010naive,golub2012how} and endogenous preferences already work with general $W$. Those literatures are this book's allies.

What structural economics adds is a single language for the three networks at once, and a structuralist agenda for it: who holds power in $W$, how $W$ forms the preferences that judge it, and how $W$ can be changed. In that sense it \textit{removes a restriction}: the anonymous benchmark restricts $W$ to rank-one matrices $\ones\mathbf{s}'$, a set of measure zero among row-stochastic matrices, and structural economics allows $W$ to be anywhere.

The geometric analogy (Euclid as the flat special case of Riemann) holds only in part. The benchmark is not identified by a single number: $\Gamma = 1$ holds for every rank-one $W = \ones\boldsymbol{\pi}'$, including the dictatorship $\ones\mathbf{e}_1'$. The ``flat'' case needs two coordinates: full integration ($\Gamma = 1$) \emph{and} equal power ($\boldsymbol{\pi} = \tfrac{1}{N}\ones$). Structure is a departure in either coordinate, and the two can move independently.
\end{intuition}


%% =====================================================================
\section{What Bolivia Taught Us}
\label{ch28:sec:bolivia}
%% =====================================================================

Bolivia appeared in every chapter --- not because Bolivia is the only developing country, but because it is a demanding test case. The data confirm some of the pathologies the book expected (an export enclave, low and falling complexity), contradict others (low integration, concentrated production power), and show equity gains that the growth-equity conjecture did not allow for.

\begin{example}[Bolivia: The Structural Portrait, Complete]\label{ch28:ex:final}
\begin{center}
\begin{tabular}{lccl}
\toprule
\textbf{Instrument} & \textbf{Bolivia} & \textbf{Benchmark} & \textbf{Chapter} \\
\midrule
\multicolumn{4}{l}{\emph{Sourced}} \\
$\Gamma(W^P)$, 12 sectors, 2019 & 0.467 & 1.00 & 17 \\
\quad like-for-like (others: 0.463--0.649) & 0.475 & & 17, 25 \\
$G(\boldsymbol{\pi}^P)$ (KOR 0.658, DEU 0.611) & 0.445 & 0.00 & 17 \\
$\pi$ of hydrocarbons in $W^P$ & 0.039 & & 17 \\
Hydrocarbons' output exported & 70\% & & 17, 20 \\
ECI, 2006 $\to$ 2019 & $-1.21 \to -1.30$ & $> 0$ & 16 \\
$\pi_G$ (global trade $W^P$) & 0.00053 & $1/N_G$ & 16 \\
$v_\Gamma$ (1990--2013) & $-0.036$/decade & & 23 \\
Gini, 2005 $\to$ 2019 & $0.600 \to 0.426$ & & 17, 26 \\
Extreme poverty, 2005 $\to$ 2018 & $38.2\% \to 15.2\%$ & & 17 \\
\midrule
\multicolumn{4}{l}{\emph{Illustrative} \illustrative} \\
$\sigma_B$ & 0.62 & 0.00 & 17 \\
$N_{\text{eff}}(W^I)$, four ecosystems & 3.4 & 4 & 17 \\
$N_{\text{eff}}(W^U)$, four groups & 3.85 & 4 & 22 \\
$\hat{\Psi}$ (self-employment) & 1.56 & --- & 17 \\
Truncation $\tau$ & 0.98 & 0.00 & 25 \\
\midrule
\multicolumn{4}{p{0.8\textwidth}}{Power ratio per capita: indigenous woman (Potos\'{i}) to German engineer $\approx$ 1 : 467 (Chapter~19, illustrative)} \\
\multicolumn{4}{p{0.8\textwidth}}{SCI (1990--2019): undefined, because Bolivia's $\Gamma$ starts above $\Gamma^{\text{ref}} = 0.50$ (Chapter~23)} \\
\multicolumn{4}{p{0.8\textwidth}}{Truncation share of the integration deficit: not identified at 12 sectors (Chapter~25)} \\
\multicolumn{4}{p{0.8\textwidth}}{2005--2019: large falls in poverty and inequality; aggregate $\Gamma$ flat to falling; ECI falling} \\
\bottomrule
\end{tabular}
\par\footnotesize Sources: $\Gamma$, $G(\boldsymbol{\pi}^P)$ and $\pi$, domestic input-share $W^P$, basic prices, diagonal kept, from the INE supply-use tables (base 2017 for 2019; base 1990 for $v_\Gamma$) \citep{ine2026sut} and, for comparators, OECD ICIO 2021 and 2023 \citep{oecd2021icio,oecd2023icio}, author's calculations (data/compute/compute\_bolivia\_io.py, compute\_oecd\_io.py, compute\_velocity.py); the like-for-like row uses the ISIC-based 12-sector aggregation, as do the comparators (Korea, Germany, the United States, Chile, China, Mexico). ECI: Growth Lab, HS92 \citep{growthlab2026atlas}. $\pi_G$: Growth Lab bilateral trade, 230 economies, author's calculations (data/compute/compute\_global\_wp.py). Gini of per-capita household income and extreme poverty (old lines; 13.2\% in 2019 on the updated lines): INE, Encuesta de Hogares, cuadros 3.06.01.02--03.
\end{center}

This table is the structural health record of an economy. In this edition its rows on the production network, complexity, global trade and distribution are sourced, and the chapter scripts in \texttt{code/} read them from \texttt{data/processed}; the between-block, information, utility and formation rows remain illustrative placeholders. Read together, the sourced rows tell a different story from the one this book set out to tell. Bolivia's domestic production network is not a fragmented outlier: its $\Gamma$ is inside the comparator range, and $\Gamma$ fell in Bolivia (1990--2013) as it did in Korea (1995--2019), so aggregate $\Gamma$ does not separate development paths. Its production power is \emph{less} concentrated than Korea's or Germany's; in industrial economies $\boldsymbol{\pi}(W^P)$ concentrates on manufacturing and business-service hubs, which is industrial depth, not oppression. What is distinctive is the enclave, measured by $\pi$ and export orientation; complexity, low and falling; and, against the book's expectations, a large fall in inequality and poverty.
\end{example}

The numbers are not the ones an earlier draft expected, and that is the point of computing them. They are \textit{specific}. The structural diagnosis does not say ``Bolivia is poor.'' It says: ``Bolivia's domestic production network is about as integrated as Korea's at the 12-sector level ($\Gamma = 0.475$ against 0.511), so aggregate integration is not its binding constraint, and the truncation conjecture gets no support at that level (Chapter~25). Its export sector is an enclave: hydrocarbons have $\pi = 0.039$ in the domestic network, export 70\% of their output, buy 0.225 of domestic inputs per unit of output and pay 6\% of their value added to employees; the rent is captured fiscally (IDH), not through linkages. Its complexity is low and falling (ECI $-1.30$, 131st). Its inequality nonetheless fell, from a Gini of 0.600 to 0.426 over 2005--2019, which Chapter~26's conjecture does not allow for. The information network has 3.4 effective voices and the formation channel produces self-confirming preferences ($\Psi = 1.56$), both illustrative until measured. The proposed intervention sequence is information $\to$ education $\to$ infrastructure $\to$ industrial policy, aimed at complexity and at the enclave's linkages, at an illustrative cost of \$17B over 30 years (about 1.2\% of 2019 GDP a year). The political precondition is a coalition shift enabled by Phase~1's $W^I$ reforms.''

Precision is the gift structural economics offers. The standard model says: ``get the prices right.'' The structuralist tradition says: ``change the structure.'' This book says: ``here is the structure, here are its eigenvalues and its distribution of power, here is what we conjecture is wrong, here is where the data have already shown us wrong, and here is a specification for fixing it that the data can check.''


%% =====================================================================
\section{What the Book Does Not Do}
\label{ch28:sec:limits}
%% =====================================================================

Intellectual honesty requires stating the limits.

\textbf{1. The model is static in important ways.} The co-evolutionary system (Chapter~13) is a difference equation, not a continuous-time model with stochastic shocks, financial frictions, or expectations. The formation model (Chapter~10) has a single preference parameter per agent. Reality has many. The model captures the structural skeleton; it does not capture the flesh.

\textbf{2. The empirical estimates are rough.} $\Gamma$, $N_{\text{eff}}$, and $\Psi$ are estimated from imperfect proxies (I-O tables that are updated every 5--10 years, media surveys with small samples, cohort-year decompositions that rely on strong identifying assumptions). The structural health dashboard is a first approximation. The values for the production network, complexity, trade and distribution of Bolivia and its comparators are now sourced (INE, OECD ICIO, the Growth Lab, the World Bank, ILOSTAT; \texttt{data/README.md}), and $\Gamma$ itself depends on aggregation and normalization (Bolivia's 2019 value ranges over 0.33--0.69). The between-block, information and utility values remain illustrative placeholders, internally consistent but not yet estimated. The quantitative claims (thresholds such as $\Gamma^*$ and $N^*_{\text{eff}}$, shares such as the truncation penalty) stand or fall with the estimates, and the first estimates have already overturned one premise: Bolivia is not a low-$\Gamma$ outlier.

\textbf{3. The political economy is underdeveloped.} Chapter~14's coalition model is a one-shot game. Real political dynamics involve repeated interaction, institutions, ideology, identity, and contingency. The hegemonic cycle (Chapter~24) is a conjecture, not a fully specified model. Reading $\boldsymbol{\pi}$ as bargaining power is an assumption; network-bargaining theory \citep{manea2011bargaining,abreu2012bargaining} shows that bargaining payoffs need not equal eigenvector centrality. The book identifies the political constraint but does not solve it.

\textbf{4. Culture is underspecified.} The formation model reduces culture to a preference parameter $\theta$ formed through $W^I$. Real culture involves narratives, identities, rituals, languages, cosmologies --- things that do not reduce to a scalar. The model captures the \textit{mechanism} of cultural transmission (network-mediated signal processing) but not the \textit{content}.

\textbf{5. The book is about structure, not about people.} Every example features an ``indigenous woman in Potos\'{i}'' or a ``gas executive in Santa Cruz'' as a node in a matrix. Real people are not nodes. They have names, histories, hopes, contradictions, and dignity that no eigenvalue can capture. The structural portrait is necessary but not sufficient. The engineer who designs the bridge still needs to talk to the people who will cross it.

\textbf{6. Many central claims are conjectures.} The book proves what can be proved: the properties of $\boldsymbol{\pi}$ and $\Gamma$, the independence of the two, the Dobrushin and Hajnal results for time-varying $W$, the saddle-node of counter-hegemony, the bridge bound on $\Gamma$, the echo-chamber result under engagement ranking. Several of its most important claims --- that truncation lowers integration (Chapter~25), that growth with equity needs integration, complexity and information diversity together (Chapter~26), that hegemonies cycle (Chapter~24) --- are conjectures. The first gets no support at the 12-sector level. Korea's early growth with equity under controlled media strains the second, and Bolivia's 2005--2019 fall in inequality with a negative, falling ECI contradicts it as a claim about equity gains. The mathematics says what would follow if they hold; the data must decide whether they do.

\textbf{7. Structure has (at least) two dimensions.} An earlier draft of this book used the spectral gap $\Gamma$ as a single measure of ``distance from the textbook.'' It is not: $\Gamma$ measures integration, and power is measured by $\boldsymbol{\pi}$; any degree of concentration is compatible with any spectral gap. No single number summarizes $W$, and a dashboard that reports only $\Gamma$ will score a winner-take-all network as healthy. Nor does aggregate $\Gamma$ separate development paths: at 12 sectors Bolivia, Korea, Germany and the United States all have $\Gamma$ between 0.46 and 0.55, and it fell in both Bolivia and Korea. Power concentration needs the same care: in a production network a high $G(\boldsymbol{\pi}^P)$ can mark industrial hubs, and Bolivia's (0.445) is lower than Korea's or Germany's. The distributional reading of $\boldsymbol{\pi}$ belongs to bargaining, under Assumption~7.A.


%% =====================================================================
\section{The Structuralist Research Program}
\label{ch28:sec:program}
%% =====================================================================

This book opens a research program; it does not close one. Each limit above is also a question. Ten for the next generation:

\begin{enumerate}
\item \textbf{Estimation} (limit 2): Finish replacing illustrative numbers with sourced estimates computed by public code. The production, complexity and distribution values are done; $\sigma_B$ (household microdata), $W^I$ (audience data), $W^U$ and the product-space quantities ($\tau$, relatedness) are not. Then estimate $W^P$, $W^I$, $W^U$ at high frequency and fine spatial resolution, using big data, satellite imagery, and digital trace data, with bootstrap error bands.

\item \textbf{Identification of $\Psi$}: Better designs for separating enforcement from formation. Randomized $W^I$ interventions (following \citealp{bernard2014future}) in multiple countries.

\item \textbf{Testing the conjectures} (limit 6): Is the capital-goods core a bridge in the product space, so that truncation lowers integration? The 12-sector data cannot tell; the test needs INE's 39 industries against comparable detail, or product-level RCA data. Does sustained growth with equity require $\Gamma > \Gamma^*$, ECI $> 0$ and $\Neff > N^*_{\text{eff}}$, and where are the thresholds? Bolivia's 2005--2019 record says equity gains do not; whether they last without structural change is the open question. What is the minimum viable capital-goods cluster for a small economy, and can it be shared regionally (a Mercosur or Andean Community program)?

\item \textbf{Algorithmic $W^I$} (limit 7): Real-time monitoring of both $\lambda_k$ (fragmentation) and $N_{\text{eff}}$ (concentration) under algorithmic curation, including messaging platforms. Experimental evaluation of diversity-mandated recommendation systems in the Global South.

\item \textbf{Climate and structure}: How does climate change rewire $W^P$ (through agricultural disruption, migration, infrastructure damage)? What are the structural effects of green industrial policy?

\item \textbf{Gender formation}: Randomized evaluations of $W^I$ interventions (role models, media campaigns, mentor programs) with long-run follow-up across generations.

\item \textbf{Multi-country structural trajectories}: Build the dashboard for every country, every year, and track the global structural trajectory. A ``structural world development indicators'' project.

\item \textbf{Political economy of operator sequencing} (limit 3): Under what political conditions is the $\Psi$-ordered sequence feasible? When does the information operator create the coalition for the structural operator? What microfoundation, if any, makes $\boldsymbol{\pi}$ a bargaining weight? What slow feedback, if any, makes hegemonies cycle?

\item \textbf{Endogenous institutions}: How do institutions (property rights, courts, regulations) emerge from the co-evolutionary dynamics? The institutional operator (Chapter~15) is exogenous in this book; it should be endogenous.

\item \textbf{The general equilibrium of structure} (limits 1, 4): Structuralist CGE models exist \citep[\textit{Socially Relevant Policy Analysis: Structuralist Computable General Equilibrium Models for the Developing World}]{taylor1990socially}. What does not yet exist is one with all three networks, the formation channel, time-varying $W$ and political equilibrium --- and with culture represented by more than a scalar.
\end{enumerate}


%% =====================================================================
\section{The Last Equation}
\label{ch28:sec:last}
%% =====================================================================

The book began with one equation (Chapter~1), the equal-endowment benchmark:
\begin{equation}
\Wce = \frac{1}{N}\ones\ones'.
\end{equation}

It ends with its generalization:
\begin{equation}\label{ch28:eq:last}
W = \Wce + \Delta(W).
\end{equation}

Every economy is the benchmark plus a departure. The departure is the structure, and it has two independent coordinates: how integrated the network is ($\Gamma$) and how power is distributed in it ($\boldsymbol{\pi}$). Together with the primitives of each model, the structure shapes allocation ($\x^*$), beliefs ($\mathbf{b}^*$), preferences ($\boldsymbol{\theta}^*$), distributional welfare ($L$), dynamics ($\Gamma_t$, $\tau_1$), politics ($S^*$), and the formation of the next generation's preferences ($\boldsymbol{\theta}^*_{t+1}$).

The benchmark model abstracts from the second term. This book studies it.

\begin{keyinsight}
Structural economics is not a school of thought. It is not a political position. It is not a regional tradition. It is the mathematics of social structure applied to economic life.

The weight matrix $W$ is the object. The Perron eigenvector is the power. The spectral gap is the integration. The formation model is the dynamics. The operator design is the engineering. The coalition game is the politics. The dashboard is the diagnostic. The trajectory is the history.

One matrix, two coordinates, and the primitives that turn them into outcomes. What follows from them is partly proved, partly conjectured, and mostly still to be measured. That is the research programme.
\end{keyinsight}

\vspace{2em}

\begin{center}
\rule{5cm}{0.4pt}
\end{center}


%% =====================================================================
\newpage
\section*{Structuralist Case Study: Furtado's Return}
\addcontentsline{toc}{section}{Structuralist Case Study: Furtado (1961)}
%% =====================================================================

Celso Furtado (1920--2004) appears three times in this book. In Chapter~2 (underdevelopment as a stable fixed point of $W$), in Chapter~9 (cultural dependency as endogenous preferences), and here, at the end.

Furtado was one of the greatest economic thinkers Latin America produced. He was trained in Paris (doctorate at the University of Paris, 1948), worked at CEPAL under Prebisch (1949--1957), led the working group on the Northeast (GTDN, 1958) and then directed SUDENE, the Northeast Brazil development agency, from its creation in 1959 until 1964, served as Brazil's Minister of Planning (1962--1963), and lost his political rights after the military coup of 1964. In exile he held a chair at the Sorbonne (1965--1985) and visiting posts at Yale and Cambridge; back in Brazil he was Minister of Culture (1986--1988). He wrote more than 30 books, across five decades. His intellectual range --- from Keynesian macroeconomics to cultural theory to institutional analysis --- anticipated almost every theme in this book.

The thesis of \textit{Desenvolvimento e Subdesenvolvimento} (\citeyear{furtado1961desenvolvimento}) can now be stated in its complete form.

\textbf{Claim 1: Underdevelopment is a structure, not a stage.} \citet{rostow1960stages} had argued that all economies pass through ``stages of growth'' from traditional to modern. Furtado rejected this: underdevelopment is not pre-development. It is a specific structural configuration produced by the \textit{interaction} of modern and archaic elements within a single system. The modern sector does not gradually absorb the traditional sector; it coexists with it, extracting surplus from it, and reproducing the duality.

In this book: underdevelopment is a structural equilibrium (Definition~\ref{ch13:def:structural_eq}) with a modern enclave weakly tied to the domestic network (in Bolivia's data, a small $\pi$ and a high export share for hydrocarbons, rather than a small aggregate $\Gamma$), concentrated bargaining power $\boldsymbol{\pi}$ (the surplus-extraction channel), and self-reproducing dynamics. It is a fixed point, not a transient state. The fixed point exists because the co-evolutionary dynamics (proximity-dominated edge creation, self-confirming preferences, coalition protection) make it stable.

\textbf{Claim 2: The structure produces the preferences that accept it.} Furtado's deepest insight, developed in \textit{O Mito do Desenvolvimento Econ\^{o}mico} (\citeyear{furtado1974mito}): underdevelopment persists not only because of economic structure but because of cultural structure --- the patterns of consumption that the economic structure produces and that in turn reproduce it. The high-income minority of the periphery imitates the consumption patterns of the center; producing those goods requires capital-intensive technologies imported from the center; those technologies absorb little labour and concentrate income further, which sustains the imitative minority. The ``myth'' is that the center's consumption standards could be generalized to the whole population of the periphery.

In this book: this is the formation model (Chapters~9--10). Preferences are endogenous to $W^I$. The formation equation $\theta_{i,t+1} = \rho\theta_{i,t} + (1-\rho)\theta^{\circ} + \mu(\sum W^I_{ij}r_j - r^*)$ captures part of Furtado's mechanism: agents form preferences by observing $W^I$-weighted signals, and when $W^I$ is dominated by center signals (Chapter~8, epistemic Perron dominance), the formation produces center-mimicking preferences in the periphery --- Furtado's ``cultural dependency.'' What the formation equation does not capture is the link Furtado cared about most: from the composition of elite demand to technology choice and to the distribution of income. A faithful formalization would make the consumption basket, not a scalar aspiration, the formed object. The $\Psi$ ratio (Definition~\ref{ch09:def:psi}; Chapter~11 asks what data can reveal about it) measures how much of a perturbation's effect comes through the formation channel.

\textbf{Claim 3: Structural change requires changing the structure, not the prices.} Furtado argued that neither market liberalization (which assumes the structure is efficient) nor simple redistribution (which operates within the structure) can solve underdevelopment. What is needed is structural transformation: changing the production structure ($W^P$), the information structure ($W^I$), and the institutional structure --- simultaneously and deliberately.

In this book: this is Part~IV (Chapters~13--15) and the neo-structuralist synthesis (Chapter~26). The operator design framework (Chapter~15) provides the engineering manual that Furtado's diagnosis called for. The $\Psi$-ordered sequence (information first, structure second, industrial policy third) is a proposed implementation protocol. The three conditions of the growth-equity conjecture ($\Gamma > \Gamma^*$, ECI $> 0$, $N_{\text{eff}} > N^*_{\text{eff}}$) are proposed success criteria, to be tested.

\begin{center}
\begin{tabularx}{\textwidth}{LLL}
\toprule
\textbf{\citet{furtado1961desenvolvimento,furtado1974mito}} & \textbf{This book's formalization} & \textbf{Chapter} \\
\midrule
Underdevelopment is a structure & Non-anonymous $W$: an enclave (low $\pi$, high export share), concentrated bargaining power & 1--4 \\
Not a stage & Stable fixed point of co-evolutionary system & 13 \\
Modern-archaic coexistence & Export enclave at the periphery of $W^P$, Pinto's heterogeneity & 3, 17 \\
Cultural dependency & $W^I$ produces center-mimicking $\boldsymbol{\theta}^*$ (consumption link not modelled) & 9, 10 \\
Frustrated modernization & $\theta^*_{\text{center}}$ incompatible with peripheral $W^P$ & 10, 12 \\
Structural change needed & Operators on $W$, not on $\boldsymbol{\omega}$ & 15 \\
Simultaneity required & Three conditions (a conjecture), three pillars & 26 \\
The general case & $W$ is any row-stochastic matrix; two coordinates & 28 \\
\bottomrule
\end{tabularx}
\end{center}

Furtado wrote much of his work in exile, in Portuguese and French, from a chair in Paris. He died in 2004, the year before Evo Morales was elected and five years before \citeauthor{hidalgo2009building}'s \citeyear{hidalgo2009building} paper on economic complexity. He did not see the product space, the ECI, or the formation experiments that would test his intuitions. He did not use the Perron eigenvector, the spectral gap, or the DeGroot model.

But he had the thesis: underdevelopment is a structure, not a stage. This book offers a formal language in which to state that thesis, to prove some of its parts, and to put the rest to the test.


%% =====================================================================
\subsection*{Further Reading}
\addcontentsline{toc}{subsection}{Further Reading}
%% =====================================================================

\textbf{Furtado:} \citet{furtado1961desenvolvimento}, \textit{Desenvolvimento e Subdesenvolvimento}. English: \textit{Development and Underdevelopment} \citep{furtado1964development}. \citet{furtado1974mito}, \textit{O Mito do Desenvolvimento Econ\^{o}mico}. English: ``The Myth of Economic Development.'' \citet{furtado1997obra}, \textit{Obra autobiogr\'{a}fica}, 3 vols.\ (his autobiography). \citet{furtado1999longo}, \textit{O Longo Amanhecer: Reflex\~{o}es sobre a Forma\c{c}\~{a}o do Brasil} (essays). \citet{love1994economic}, ``Economic Ideas and Ideologies in Latin America since 1930'' (for Furtado's intellectual context).

\textbf{The structuralist tradition, complete:} \citet{prebisch1949desarrollo}, the manifesto. \citet{furtado1961desenvolvimento}, the thesis. \citet{pinto1970naturaleza}, the diagnosis. \citet{sunkel1971capitalismo}, the integration/disintegration. \citet{tavares1972substituicao}, the autopsy. \citet{marini1973dialectica}, the exploitation. \citet{fajnzylber1983industrializacion,fajnzylber1990industrializacion}, the truncation and the empty box. \citet{cardoso1969dependencia}, the politics. \citet{quijano2000colonialidad}, the coloniality. \citet{sen1999development}, the capabilities. \citet{cepal2012structural}, the institutional synthesis. \citet{ocampo2009growth}, the policy synthesis. \citet{hidalgo2007product} and \citet{hidalgo2009building}, the empirical measurement. \citet{taylor1983structuralist,taylor1990socially}, the structuralist macroeconomics and its CGE models.

\textbf{Networks in economics:} \citet{long1983real}; \citet{acemoglu2012network}, ``The Network Origins of Aggregate Fluctuations''; \citet{baqaee2019macroeconomic}, ``The Macroeconomic Impact of Microeconomic Shocks''; \citet{golub2010naive,golub2012how}; \citet{seneta2006nonnegative}, \textit{Non-negative Matrices and Markov Chains}.

\end{document}
